%------------------------
% Taal & codering
%------------------------
\usepackage[english]{babel}        % Nederlandse hyphenatie en typografie
\usepackage[T1]{fontenc}         % Correcte font-encoding (accenten, afbreking)
\usepackage[utf8]{inputenc}      % Bronbestand in UTF-8 (met moderne LaTeX niet strikt nodig)
\usepackage{csquotes}            % Correcte aanhalingstekens (werkt goed met biblatex)

%------------------------
% Lettertype
%   XCharter + newtxmath geven een nette, consistente set
%------------------------
\usepackage{XCharter}            % Tekstfont
\usepackage[charter]{newtxmath} % Wiskunde bij Charter (variabele symbolen)
\usepackage{microtype}           % Mooiere microtypografie

%------------------------
% Wiskunde & symbolen
%------------------------
\usepackage{amsmath,mathtools}
\usepackage{siunitx}
\sisetup{
    per-mode = symbol,
    exponent-product = \cdot,
    output-decimal-marker = {.},
    exponent-product = \cdot
}
\DeclareSIUnit\phase{ph}  % Extra eenheid voorbeeld
\DeclareSIUnit{\ohmpsquare}{\text{\ensuremath{\Omega/\square}}}
\usepackage{nicematrix}

%------------------------
% Afbeeldingen, figuren, grafieken
%------------------------
\usepackage{graphicx}
\usepackage{wrapfig}
\usepackage{float}
\usepackage{subcaption}   % Vervangt subfig
\usepackage{pgfplots}
\usepackage{matlab-prettifier}
\usepgfplotslibrary{polar}
\pgfplotsset{compat=1.18}

%------------------------
% TikZ & elektrische schema's
%------------------------
\usepackage[european,american voltages]{circuitikz}
\usetikzlibrary{
  arrows.meta,
  positioning,
  shapes.geometric,
  patterns,
  backgrounds,
  scopes
}

%------------------------
% Tabellen en lijsten
%------------------------
\usepackage{booktabs,longtable,multirow,multicol,array}
\usepackage{tabularx}
\newcolumntype{Y}{>{\centering\arraybackslash}X}
\newcolumntype{?}{!{\vrule width 2pt}} % Dikke verticale lijn indien echt nodig
\newcommand{\thickhline}{\noalign{\hrule height 2pt}}

%------------------------
% Codefragmenten
%------------------------
\usepackage{listings}
\lstset{
  basicstyle=\ttfamily\small,
  frame=single,
  columns=fullflexible,
  showstringspaces=false,
  tabsize=2,
  breaklines=true,
  language=C
}

\definecolor{codegreen}{rgb}{0,0.5,0}
\definecolor{codegray}{rgb}{0.5,0.5,0.5}
\definecolor{codepurple}{rgb}{0.58,0,0.82}
\definecolor{backcolour}{rgb}{0.97,0.97,0.97}

\lstdefinestyle{mypython}{
    language=Python,
    backgroundcolor=\color{backcolour},
    basicstyle=\ttfamily\scriptsize,
    keywordstyle=\color{blue},
    commentstyle=\color{codegreen},
    stringstyle=\color{codepurple},
    numberstyle=\tiny\color{codegray},
    numbers=left,
    numbersep=6pt,
    breaklines=true,
    breakatwhitespace=false,
    showstringspaces=false,
    showspaces=false,
    showtabs=false,
    tabsize=4,
    frame=single,
    captionpos=b,
    keepspaces=true
}

%------------------------
% Verwijzingen & hyperlinks (volgorde: eerst hyperref, dan cleveref)
%------------------------
\usepackage[hidelinks]{hyperref}
\usepackage[noabbrev,nameinlink]{cleveref}

\usepackage{attachfile2}

%------------------------
% Opmaak & randvoorwaarden
%------------------------
\usepackage{parskip}      % Geen inspringing, wel witregel tussen alinea's
\usepackage{enumitem}     % Fijnere controle lijsten
\usepackage{xcolor}
\usepackage{fancyhdr}
\usepackage{geometry}
\geometry{
  a4paper,
  left=30mm,right=30mm,top=33mm,bottom=20mm
}
\pagestyle{fancy}
\fancyhf{}
\lhead{\leftmark}
\rhead{\thepage}

%------------------------
% Bibliografie (IEEE-stijl, biber-backend)
%------------------------
\usepackage[
  backend=biber,
  bibstyle=ieee,
  citestyle=numeric-comp,
  doi=false,isbn=false,url=true
]{biblatex}
\addbibresource{references.bib}

%------------------------
% Caption-breedte en centrering
%------------------------
\usepackage[justification=centering]{caption}
\captionsetup{width=0.95\linewidth}

%------------------------
% Handige eigen omgevingen/lijsten
%------------------------
\newlist{requirements}{enumerate}{10}
\setlist[requirements]{label*=REQ-\arabic*,leftmargin=*,itemsep=0.25em}

\newlist{Nreq}{enumerate}{10}
\setlist[Nreq]{label*=NREQ-\arabic*,leftmargin=*,itemsep=0.25em}

\newlist{research}{enumerate}{10}
\setlist[research]{label*=RQ-\arabic*,leftmargin=*,itemsep=0.25em}

%------------------------
% Bestandslocaties voor afbeeldingen
%------------------------
\graphicspath{{./figures/}}

% Optioneel: euroteken (meestal overbodig met moderne fonts)
% \usepackage{eurosym}

% Optioneel: PDF inladen
% \usepackage{pdfpages}

% Optioneel: externe referenties
% \usepackage{xr}

\mathchardef\Re="023C
\mathchardef\Im="023D

\DeclareSIUnit{\dBi}{dBi}
\DeclareMathOperator{\sinc}{sinc}

\usepackage{titlesec}
\usepackage{tocloft}
\setlength{\cftsubsecnumwidth}{2.25cm}

\renewcommand{\thesubsection}{Question \arabic{subsection}}

\titleformat{\subsection}
{\normalfont\large\bfseries}
{\thesubsection}
{1cm}
{}

\renewcommand{\cftdot}{}